
\section{Conclusion}\label{sec:conclusion}
This paper has developed a structural diagnostic for stablecoin basket risk whose central object is a correlation rigidity that every smooth-surface method must forgo: within a failure regime the joint de-peg correlation matrix of a basket does not move, and across a boundary it steps. The decisive technical choice is which matrix carries the crossing law. The wall-crossing operation is additive and rank-one, and that operation is admissible on the co-movement \emph{covariance} but not on the correlation matrix: two correlation matrices both have unit diagonal, so their difference must have zero diagonal, whereas a rank-one update supported on a block of two or more coins contributes a strictly positive amount to the diagonal. Stating the law on the covariance and passing to correlations through the normalising map removes the obstruction exactly, because that map \emph{assigns} the diagonal rather than computing it. What the choice buys is that the crossing can be stated at all: on the covariance it is admissible, and what then follows is that rigidity on chambers is inherited by the correlation matrix, that the induced jump has zero diagonal, and that the part of it surviving normalisation is its \emph{sign structure and support} rather than its rank or its scale. \emph{Contagion}, in this vocabulary, is not a market-wide rise in correlation but a structured one: the block signature is what a channel's failure does to the co-movement of the coins that channel touches, and its confinement to those coins is what makes it measurable rather than merely observed. A peg-preserving parameter space is stratified by eight stress boundaries, one per structural failure channel, and the block at each is prescribed ex ante by the wall taxonomy rather than fitted.

The paper's logical status is worth stating plainly. The three load-bearing properties---within-chamber rigidity, the rank-one form of the crossing, and its support on the wall-affected block---are \emph{assumed} (\S\ref{sec:what-is-assumed}), not derived. What is proved is everything downstream: that rigidity of the covariance gives rigidity of the correlation matrix exactly, with the scale cancelled; that the induced correlation jump has zero diagonal; that its image is spread beyond the block in a signed three-tier pattern; and that no smooth surface can reproduce it. What is tested is whether those consequences describe the 2020--2025 record. Splitting the paper this way is what lets a reader who doubts the assumptions still use the results, and a reader who accepts them still be held to the data.

The empirical design follows the theory's corrected shape rather than forcing the theory into a shape it cannot take. Because only support and sign survive normalisation, the test is a pair-level gradient: classify each pair of coins by how many lie in the wall-affected block and ask whether the acute-window change in correlation is ordered accordingly, with the block-versus-rest contrast differencing out the crisis-wide co-movement that would otherwise contaminate every layer. On the nine-event 2020--2025 public panel, the primary event (USDC--SVB, $W_4$-\texttt{bank}) gives a clean three-step staircase, a block-versus-rest contrast of $\GUSDCSVBDID$, and a coin-level permutation $p = \GUSDCSVBPCOIN$; the result survives a full $3\times3$ window grid, and the pair $(\text{DAI},\text{FRAX})$---which contains neither the primary victim nor any coin with direct custody exposure---jumps by $\GUSDCSVBPURE$, the largest of all pairs, defeating the reading that the gradient is a victim effect. The other two multi-coin-block events supply a single BLOCK pair each and are reported alongside the primary event in Table~\ref{tab:gradient}.

Three contributions are offered. First, a wall-stratified taxonomy of eight structurally distinct de-peg mechanisms, and the claim---stated as an assumption so that it can be contested---that the joint co-movement structure carries an ex-ante prescribed block signature across each wall, intrinsic to the wall type rather than to the estimator used. Second, a rigidity-and-crossing framework carried to its testable consequences: a proof that correlation rigidity follows from covariance rigidity with the scale cancelled, a proof that the induced correlation jump has zero diagonal, and a precise account of what any locally continuous, smooth-surface recovery scheme must forgo on a multi-coin-block event. The signature is computable by a supervisor from public price data with no distributional assumption and no issuer cooperation. Third, a pair-level empirical test that rejects the no-block null on the primary event, with the other two multi-coin-block events reported alongside it rather than pooled with it.

Two implications follow. Basket-margin frameworks, proof-of-reserves audits, and index-perpetual clearinghouses that price joint de-peg risk on a smooth surface systematically understate the tail in the multi-coin-block regime by a signature the theory identifies in sign and support, and that no smooth surface can place. The operational remedy is deliberately small: a \emph{block-alignment statistic}---the excess acute-window co-movement of block pairs over all other pairs, which is exactly $\widehat{\mathrm{DiD}}$ of \eqref{eq:did}---computed on a public correlation submatrix and paired with a regime-dependent margin multiplier. It requires neither proprietary data nor issuer cooperation, and it degrades gracefully: on the events where no block is defined it returns nothing rather than a spurious number. On the research side, the construction generalizes to any basket stratified by structural failure modes (cross-margined perpetuals, tokenized money-market funds, bank-consortium reserve pools); moving from a support test to distributional inference requires more multi-coin-block events, which the 2020--2025 record does not contain. The eight-wall taxonomy leaves cross-chain bridge failures and stablecoin-native regulatory perimeters as natural next entries.

Basket risk is a staircase, not a slope---identifiability jumps at walls rather than sliding between them. Pricing the slope is pricing the wrong risk; identifying the walls, and refusing to price what they do not identify, is the discipline the structure enforces. The block signature is the height of each step, and the rank-one residue is what lifts one step to the next. What the structure delivers is a coordinate system in which margin, audit, and supervision can be written down without pretending the walls are not there---a matter of immediately practical significance for cryptocurrency risk measurement.

\appendix
\titleformat{\section}[block]{\normalfont\Large\bfseries}{Appendix \thesection:}{0.8em}{}
\section{Data Construction}\label{app:data}
The \NPANEL\ stablecoins are drawn from the DeFiLlama public archive of daily peg-price series over 2020--2025, one CSV per stablecoin. The universe was constructed by taking all stablecoins with more than 300 daily observations, non-trivial trading volume, on public exchanges, and a declared 1:1 peg to either USD, EUR, gold, or another fiat/asset reference; series that failed the coverage screen are not carried into the panel. The panel partitions into four tiers:
\begin{itemize}[leftmargin=1.8em]
\item Tier 1 (fiat-backed, 11 coins): USDT, USDC, BUSD, TUSD, USDP, GUSD, FDUSD, PYUSD, EURS, GYEN, XSGD;
\item Tier 2 (crypto-collateralized, 10 coins): DAI, LUSD, GHO, crvUSD, MIM, USDD, RAI, VAI, USDe, HUSD;
\item Tier 3 (algorithmic, 4 coins): UST, FEI, FRAX, NxUSD;
\item Tier 4 (RWA/other, 3 coins): EUROC, EURoe, CNHT.
\end{itemize}
A given event window uses only the subset with at least five usable paired observations in both its pre- and acute windows; on the primary event that subset has 18 coins, which is the $n$ entering the $\binom{18}{3}=816$ permutation of \S\ref{sec:gradient}.

\paragraph{Event dating.} Each of the nine events in Table~\ref{tab:events} is dated by the primary structural failure identified in the regulator and post-mortem literature: DAI Black Thursday (MakerDAO retrospective, Chainlink oracle logs); UST/LUNA \citep{LiuLiSprenger2023}; USDT stress (exchange net-outflow filings); FTX contagion (FTX bankruptcy filing 11 November 2022 following Alameda balance-sheet leak; BIS Quarterly Review March 2023 as post-mortem literature; $W_4$-\texttt{cex} block assignment $\{\text{USDD, MIM}\}$ from identifiable Alameda/FTX counterparty exposure at the coin-issuer level); BUSD (NY DFS enforcement action Feb 13 2023); USDC-SVB (SVB closure order Mar 10 2023, Circle Mar 11 disclosure); FDUSD 2024 (First Digital Trust public statement Nov 8 2024). Event windows are $\pm 60$ days for time-series plots and $[-5, +5]$ days for acute-stress correlation estimation.

\paragraph{Data provenance.} No proprietary, paywalled, or issuer-internal data is used. Level-1 (real) sources: DeFiLlama daily peg-price archive (all \NPANEL\ stablecoins), FRED macro series (VIXCLS, DGS10, T10Y2Y, BAA10Y, AAA10Y, BAAFF, AAAFF, DFF), Yahoo Finance closes (BTC-USD, ETH-USD, VIX, TNX, GLD), and Etherscan on-chain records (5 stablecoin token supplies, 3 Chainlink price-feed proxy addresses, 102 MakerDAO PSM-USDC transactions around SVB). The 102 PSM-USDC transactions are retained as auditable event-anchor evidence for the $W_1$ transmission identification of the DAI/USDC-SVB co-de-peg (a PSM flow spike on 11--12 March 2023 corroborates the wall-block assignment) and are not used as a quantitative input. Level-2 (proxy): Chainlink oracle history is proxied by DeFiLlama peg-price series with the Chainlink proxy contract addresses recorded for future verification; CoinGecko daily USD close cross-check against DeFiLlama on 99.2\% of coin-days within 3 bp tolerance. Level-3 (simulated/grounded projection): \emph{none}---every empirical number in \S\ref{sec:empirical} derives from Level-1 or Level-2 series.

\paragraph{Snapshot freeze.} The DeFiLlama archive, FRED series, Yahoo Finance quotes, and CoinGecko quotes are all live mutable endpoints. The reproducibility bundle records, for every one of the 86 raw CSV files, the source URL, retrieval timestamp (UTC), file size, SHA-256 checksum, and date range; empirical reproducibility is guaranteed against these fingerprints. Data not pulled at this submission (Dune queries, Chainlink full history) is enumerated as ``partial'' in the bundle with checkpoint state persisted; no main-text claim depends on the partial sources.

\section{Reproducibility Archive}\label{app:repro}
The submission bundle contains the \LaTeX\ source, all figure-generation and audit scripts, the raw data archive, a SHA-256 manifest, and a submission checklist. All scripts run under Python 3.11 with pinned dependencies (requirements.txt); full pipeline (excluding Chainlink full-history pull, which is checkpointed and background-runnable) completes in under three minutes on a laptop. Data-provenance details are shipped with the bundle.

\section{Out-of-Sample and Window Robustness}\label{app:oos}
This appendix reports the checks behind \S\ref{sec:robust}: the window grid for the pair-level gradient, the status of the amplitude parameter $\sigma_W$, and the near-terminal UST/LUNA regime.

\subsection{The window grid}\label{app:oos-grid}
The window specification enters the test only through the two correlation estimates $\rho^-$ and $\rho^+$. Table~\ref{tab:robust} reports the block-versus-rest contrast \eqref{eq:did} for all three multi-coin-block events on a $3\times3$ grid: pre-window length $30$, $60$ or $90$ calendar days, crossed with an acute half-width of $\pm3$, $\pm5$ or $\pm10$ days. The canonical specification of the main text is $60$\,d and $\pm5$\,d, shown in bold; the script that produces the table re-derives the canonical row exactly, which is the check that the grid is comparable across specifications. A dagger marks a specification in which the full ordering \eqref{eq:gradient-prediction} fails.

\begin{table}[!ht]
\centering
\caption{Window grid for the pair-level gradient. Entries are the block-versus-rest contrast $\widehat{\mathrm{DiD}}$ of \eqref{eq:did}. Bold marks the canonical specification used in the main text. $\dagger$ marks a specification under which the criterion \eqref{eq:gradient-prediction} does not hold. The grid varies only the correlation windows; block assignments are fixed ex ante and are not re-estimated at any specification.}
\label{tab:robust}
\begin{tabular}{@{}ccrrr@{}}
\toprule
Pre window & Acute half-width & USDC--SVB & FTX & FDUSD \\
\midrule
30\,d & $\pm$3\,d & $+0.642$ & $+0.832$ & $+0.039$ \\
30\,d & $\pm$5\,d & $+0.674$ & $+0.195$ & $+0.036$ \\
30\,d & $\pm$10\,d & $+0.789$ & $-0.472^{\dagger}$ & $+0.039$ \\
60\,d & $\pm$3\,d & $+0.632$ & $+0.847$ & $-0.039^{\dagger}$ \\
60\,d & $\pm$5\,d & $\mathbf{+0.687}$ & $\mathbf{+0.513}$ & $\mathbf{-0.246^{\dagger}}$ \\
60\,d & $\pm$10\,d & $+0.687$ & $+0.641$ & $-0.309^{\dagger}$ \\
90\,d & $\pm$3\,d & $+0.582$ & $+0.862$ & $-0.048^{\dagger}$ \\
90\,d & $\pm$5\,d & $+0.636$ & $+0.562$ & $-0.188^{\dagger}$ \\
90\,d & $\pm$10\,d & $+0.610$ & $+0.657$ & $-0.259^{\dagger}$ \\
\bottomrule
\end{tabular}

\end{table}

\paragraph{Reading the grid.} (i) The USDC--SVB criterion holds in all nine specifications, with $\widehat{\mathrm{DiD}}$ between $+0.58$ and $+0.79$, so the primary result is not a window artifact. (ii) On FTX it holds in eight of the nine; the exception is the 30-day pre-window with a $\pm10$-day acute window, where the block-versus-rest contrast itself turns negative. (iii) On FDUSD it holds only under a 30-day pre-window and fails under the 60- and 90-day windows, where the contrast is negative.

\paragraph{Why the grid is not a search over specifications.} The canonical $60$\,d\,/\,$\pm5$\,d window was fixed by the panel-construction script before the gradient test was designed, and it is the window used for every event in the paper. The grid is reported because a reader cannot otherwise judge how much of the result is a window choice; it is not used to select among specifications.

\paragraph{Expanding-window check.} Four of the nine events (USDC--SVB, USDe 2024, FDUSD 2024, GHO 2024) fall inside an expanding-window out-of-sample scheme; DAI Black Thursday precedes the daily-archive coverage, and UST, USDT, FTX and BUSD precede the training cut-off. Across expanding windows the wall-type assignment of Table~\ref{tab:events} and the resulting block membership are unchanged, which is what the ex-ante coding discipline of \S\ref{sec:ninewalls} requires; the empirical comparison itself is the one reported in \S\ref{sec:empirical}, not a re-run of it.

\subsection{Status of the amplitude parameter $\sigma_W$}\label{app:oos-sigma}
Assumption~\ref{ass:wallcrossing} carries a positive scale $\sigma_W$ that the paper leaves free. Because the normalising map of \eqref{eq:normalised-corr} is a nonlinear congruence, it destroys the rank-one form of the jump (Remark~\ref{rem:small-cases}), so the properties of the jump that survive into correlation space are its sign structure and its support, not its scale. A test built on those is invariant to $\sigma_W$ at leading order in $a$, which is the order at which the empirical criteria of \S\ref{sec:empirical} are stated. Fixing $\sigma_W$ would matter only for a test of magnitude, and that requires more multi-coin-block events than the 2020--2025 record contains.

\subsection{The UST/LUNA near-terminal regime}\label{app:oos-ust}
Corollary~\ref{cor:copula-fail} linearises the wall-crossing residue in the wall-distance parameter $\epsilon$ and is valid while peg deviations remain in the small-signal regime. Near-complete de-peg is not such an excursion: UST reaches roughly a $99\%$ drawdown at terminal on 09--12 May 2022. UST is also algorithmic-reflexive in a way that no collateral composition in Definition~\ref{def:stable} represents, so of the nine events it is the least well described by the construction of \S\ref{sec:moduli}. The event is retained in Table~\ref{tab:events} as a documented boundary of the theory rather than as a target of its prediction. Independently of this, its wall-affected block is a single coin, so it could not carry the pair-level test even if the small-signal condition held---which is the same scope statement made in \S\ref{sec:single}.

\bibliographystyle{apalike}
\bibliography{refs}

%%------------------------------------------------------------
%% DECLARATIONS — collected by the publisher's submission system
%% (CRediT roles, funding, competing interests, data availability,
%% ethics, generative-AI use, acknowledgements). Removed from the
%% manuscript body at the author's request so the file stays clean;
%% the corresponding statements are entered as form fields at submission.
%%------------------------------------------------------------

\end{document}
